\subsubsection{Expenses, Expenditures and Cost}

\textbf{Cost} can be defined "as the expenditure incurred on various inputs such as land, labor, capital, and entrepreneurship to produce goods and services".\footcite{kin2025cost} Simply put, cost is \textit{the amount of resources spent (the input) on generating a good or service}.

\textbf{Expenses} are the costs required to generate revenue and net profit for a business. Another way of framing this is, expenses are how much a business cost offers a business's revenue. There are different types of expenses, such as operating expenses (rent, utilities, etc.), costs of goods sold, administrative expenses, and selling expenses.\footcite{thomas_2026_expense}

\textbf{Expenditures} include all outflows of money, and specifically deal with the purchase of goods or services. The goal for an expenditure is to benefit the business. There are different types of expenditures: capital, and revenue expenditures. Capital expenditures (CapEx) are expenditures that are founded within long-term goods, or services, where revenue expenditures (OpEx) encompass more short-term, day-to-day costs.

\textbf{Cost structures} are the categorization of expenses a business takes on while operating. Costs such as money, real, fixed, variable, explicit, implicit, private, social, total, average, marginal, short-run, long-run, intangible, and tangible costs are all examples of cost structures (a mouthful, right?).



\underline{\textbf{Money and Real Cost}}

\textbf{Money costs} are the monetary expenditures to produce a product or service (e.g. wages, rent).

\textbf{Real costs} are the pain and effort of all resources used in production, and captures both mental and physical efforts of labor.\footcite{kin2025cost}



\underline{\textbf{Fixed and Variable Costs (Production Costs)}}

\textbf{Fixed costs} are costs that stay consistent regardless of the output produced (e.g. rent, salaries, software subscriptions). It's important to note, sunk costs can be considered as fixed costs, but not all fixed costs are to be considered as a sunk cost.\footcite{nickolas_2025_variable}

\textbf{Variable costs} are costs that vary (either increase or decrease) based on the output produced (e.g. raw materials, commission, shipping, labor). Variable costs are also highly industry specific, as product output cannot be compared along distinct industries (for instance, you cannot compare the variable costs with the production of phones to that of the production of t-shirts).\footcite{nickolas_2025_variable}

Calculating variable costs is simple; it is the quantity of output multiplied by the variable cost per unit of output.\footcite{nickolas_2025_variable} Suppose it costs \$5 dollars to produce one t-shirt, and there are 100 units of t-shirts produced. The variable cost to produce 100 units of t-shirts would be \$500.

\textbf{Semi-variable costs} are costs that involve a mixture of fixed costs, and variable costs. With semi-variable costs, the fixed portion remains stagnant, while the variable cost varies based on the output produced.\footcite{tuovila_2022_semivariable} Suppose we have a firm that produced blenders. If the firm wants to produce more blenders, we need more workers (which is a variable cost, as it depends on the output produced). However, even if I didn't have any output produced, I still need workers to look after my non-existent production, which is a fixed cost.



\underline{\textbf{Explicit and Implicit Costs}}

\textbf{Explicit costs} are transactions made out-of-pocket that have a measurable cost to a firm (e.g. wages, rent, utilities).\footcite{pettinger_2020_the} These costs can either be fixed, or variable.

\textbf{Implicit costs} are the opportunity cost of resources already owned, and do not show up on balance sheets. The best way to explain implicit costs is by looking at a non-profit company that produces a tangible product, that could be commercialized for profit. If that non-profit company uses its machinery to produce said product, it loses out on the next-best alternative which is revenue from commercializing the product.

\underline{\textbf{Private and Social Costs}}

\textbf{Private Costs} are costs directly by the business or individual in producing a good or service.

\textbf{Social Costs} are costs imposed on third parties not directly involved in the activity or society as a whole (e.g. pollution, congestion), where there are positive (creates benefits for others without being reimbursed) and negative externalities (imposes a cost on someone, without reimbursing them) as it relates to social cost.

\underline{\textbf{Total, Average, and Marginal Costs}}

\textbf{Total Cost (TC)} is the full cost of producing a given level of output; TC = Fixed Costs + Variable Costs.

\textbf{Average Cost (AC)} is cost per unit of output produced; AC = TC $\div$ Quantity Produced.

\textbf{Marginal Cost (MC)} is the extra cost you face from doing one more unit of something. For example, if the total cost of 6 units is \$1000, and the total cost of 7 units is \$1900, then the marginal cost of the 7th unit is \$900.

% \underline{\textbf{Intangible and Tangible Costs}}

\subsubsection{Value}

\textbf{Value} can be defined as the set of benefits a product or service provides a customer based on their needs.  